\documentclass[aspectratio=169,11pt]{beamer}
\usetheme{Madrid}
\usecolortheme{default}
\usepackage[T1]{fontenc}
\usepackage{amsmath}
\usepackage{booktabs}
\usepackage{graphicx}
\title[Beamer Lab]{Introduction to Beamer}
\subtitle{Technical Writing Laboratory}
\author{Ananya Siju}
\institute{Department of Computer Science and Engineering\\MNIT Jaipur}
\date{Academic Year 2026--2027}
\begin{document}
\begin{frame}\titlepage\end{frame}
\begin{frame}{Presentation Outline}\tableofcontents\end{frame}
\section{Introduction}
\begin{frame}{What is Beamer?}
\begin{itemize}
\item Beamer is a LaTeX document class for creating slides.
\item Slides are called \textbf{frames}.
\item A source file is compiled to a portable PDF.
\item Ideal for mathematical and technical presentations.
\end{itemize}
\end{frame}
\section{Core Features}
\begin{frame}{Frames and Structured Lists}
\begin{enumerate}
\item Declare a frame with \texttt{\\begin\{frame\}}.
\item Add a title and content.
\item Close with \texttt{\\end\{frame\}}.
\end{enumerate}
\end{frame}
\begin{frame}{Blocks and Emphasis}
\begin{block}{Standard Block}Group related ideas in a named block.\end{block}
\begin{alertblock}{Important}Keep slides concise and readable.\end{alertblock}
\begin{exampleblock}{Example}Use \texttt{\textbackslash textbf} for bold text.\end{exampleblock}
\end{frame}
\begin{frame}{Overlays: Step-by-Step Reveal}
\begin{itemize}
\item<1-> First idea appears on overlay 1.
\item<2-> Second idea appears on overlay 2.
\item<3-> Third idea appears on overlay 3.
\end{itemize}
\end{frame}
\section{Technical Content}
\begin{frame}{Mathematical Typesetting}
The quadratic formula is
\[x=\frac{-b\pm\sqrt{b^2-4ac}}{2a},\qquad a\ne 0.\]
An inline expression: $E=mc^2$.
\end{frame}
\begin{frame}{Tables in Beamer}
\begin{table}
\centering
\begin{tabular}{ll}
\toprule
Feature & Purpose\\
\midrule
Frames & Organize slides\\
Themes & Control presentation appearance\\
Overlays & Reveal content progressively\\
Blocks & Highlight related information\\
\bottomrule
\end{tabular}
\caption{Common Beamer features}
\end{table}
\end{frame}
\begin{frame}{Two-Column Layout}
\begin{columns}[T]
\begin{column}{0.48\textwidth}
\textbf{Advantages}
\begin{itemize}\item Consistent design\item Excellent equations\item Version-control friendly\end{itemize}
\end{column}
\begin{column}{0.48\textwidth}
\textbf{Considerations}
\begin{itemize}\item Requires LaTeX syntax\item Compilation needed\item Fine-grained visual design can take time\end{itemize}
\end{column}
\end{columns}
\end{frame}
\section{Workflow}
\begin{frame}{Practical Workflow}
\begin{enumerate}
\item Write a \texttt{.tex} source file.
\item Select a theme and add frames.
\item Compile using \texttt{pdflatex} or Overleaf.
\item Review the PDF and fix any errors.
\item Commit both source and final output to GitHub.
\end{enumerate}
\end{frame}
\begin{frame}{Learning Outcomes}
\begin{itemize}
\item Create slides using the Beamer class.
\item Use themes, blocks, columns, equations, and tables.
\item Understand overlay specifications.
\item Compile a reproducible PDF presentation.
\end{itemize}
\end{frame}
\begin{frame}[plain]
\centering\Huge Thank You!\\[1em]
\normalsize Technical Writing Laboratory\\MNIT Jaipur
\end{frame}
\end{document}
